\documentclass[conference]{IEEEtran}
\IEEEoverridecommandlockouts

\usepackage{cite}
\usepackage{amsmath,amssymb,amsfonts}
\usepackage{algorithmic}
\usepackage{algorithm}
\usepackage{graphicx}
\usepackage{textcomp}
\usepackage{xcolor}
\usepackage{booktabs}
\usepackage{url}
\usepackage{hyperref}
\usepackage{multirow}
\usepackage{array}
\usepackage{tabularx}
\usepackage{enumitem}

\emergencystretch=1em

\hypersetup{
    colorlinks=true,
    linkcolor=blue,
    citecolor=blue,
    urlcolor=blue
}

\def\BibTeX{{\rm B\kern-.05em{\sc i\kern-.025em b}\kern-.08em
    T\kern-.1667em\lower.7ex\hbox{E}\kern-.125emX}}

\begin{document}

\title{Predictive Early-Exit Routing for SLO-Compliant LLM Inference:\\
Empirical Compliance and EVT Tail Characterisation via Post-Hoc KV-Cached Forward Hooks}

\author{
    \IEEEauthorblockN{Nithin Palyam}
    \IEEEauthorblockA{
        \textit{Department of Computer Science and Engineering}\\
        \textit{University of South Florida}\\
        Tampa, FL, USA\\
        palyam@usf.edu
    }
}

\maketitle

% ─────────────────────────────────────────────────────────────────────────────
\begin{abstract}
Standard autoregressive LLM inference executes every transformer layer for
every generated token, yielding per-token latencies that are highly variable
and potentially unbounded as context grows.
For latency-sensitive applications such as clinical decision support,
this tail behaviour threatens Service-Level Objective (SLO) deadlines
in ways that standard profiling cannot certify away.
We address this mismatch by applying the \emph{anytime algorithm}
principle to LLM token generation: rather than waiting for all 22 transformer
layers to complete, intermediate-layer predictions provide routing signals
that the scheduler uses to make deadline-aware commit decisions at the
application level.

Two routers are implemented and evaluated on TinyLlama-1.1B-Chat on a
bare-metal \textbf{NVIDIA A100~SXM4 80\,GB} (HBM2e, CUDA~12.x):
(i)~a \emph{stateless two-pass} router that runs Layer~16 as a probe
     before optionally issuing the full pass; and
(ii)~a \emph{KV-cached single-pass} router that captures Layer-16
     activations via a post-hoc forward hook during the contiguous
     22-layer pass, eliminating KV-cache desynchronisation entirely.
At $D=45$~ms, the KV-cached router achieves P99~TPOT~$=18.26$~ms
($R=0.406$) with 0\% deadline misses across 150~TPOT samples;
the empirically observed minimum schedulable deadline is 25~ms.
Voluntary early exits account for 0.7\% of tokens at threshold $\tau=0.55$;
a trained 2-layer MLP exit head increases the voluntary exit rate to 79\%
at the cost of reduced oracle agreement; results are reported in
Section~\ref{sec:exithead}.
An EVT analysis of 500 spaced timing samples per cell (200\,ms inter-run
spacing) reveals heavy tails throughout: GEV shape $\hat{\xi} \in [0.21,\,1.18]$
across all 12 (seq\_len, exit\_layer) cells, Gumbel rejected by
Anderson-Darling at 5\% in every cell, and Ljung-Box sample independence
violated despite spacing (lag-1 ACF ${\approx}0.29$, $p{=}0.0$).
Block-maxima GEV is applied as a fallback but yields bootstrap CIs
spanning orders of magnitude, precluding formal probabilistic certification.
The paper's central contribution is therefore empirical SLO compliance:
P99~TPOT~$=18.26$~ms with 0\% deadline misses at $D=45$~ms (Section~\ref{sec:evt}).
A 500-query PubMedQA evaluation achieves 0\% deadline misses at
56.81~tok/s throughput and 54.6\% answer accuracy
(95\% CI: 48.3\%--60.8\%, $n{=}240$ scoreable tokens).
The clinical setting is used as a motivating benchmark scenario;
TinyLlama's accuracy places it below thresholds required for clinical
deployment.
\end{abstract}

\begin{IEEEkeywords}
anytime algorithms, SLO compliance, LLM inference,
early exit, KV cache, tail latency, probabilistic WCET, extreme value theory,
transformer inference, schedulability
\end{IEEEkeywords}

% ─────────────────────────────────────────────────────────────────────────────
\section{Introduction}
% ─────────────────────────────────────────────────────────────────────────────

Large language models have become foundational components in
latency-sensitive applications: clinical decision support~\cite{singhal2023large},
interactive human-computer interfaces~\cite{ouyang2022training}, and
autonomous vehicle reasoning~\cite{wen2023dilu}.
In all such contexts the system must commit a token---or produce an
answer---before a deadline, with limited tolerance for the tail latencies
that routinely arise in production deployments.
Unfortunately, standard autoregressive decoding executes all transformer
layers for every token and produces execution times that are both highly
variable and context-dependent, growing with the length of the KV cache.
There is no known closed-form bound, and empirical WCET measurements
alone provide only a single-sample guarantee.

This mismatch between standard inference and SLO requirements is the
central problem we address.
We adopt the \emph{anytime algorithm} paradigm~\cite{boddy1994deliberation,
zilberstein1996using}: a computation that can be safely interrupted at
any point and returns a valid (if approximate) result.
For transformer inference, the natural interruption point is at an
intermediate layer; a confidence signal derived from that layer's logits
then allows the router to decide whether the intermediate output is
sufficient for commitment, without issuing any additional forward calls.
The present implementation tests a single intermediate exit point (Layer~16);
a dense sweep of layers L5--L20 was subsequently completed
(Section~\ref{sec:ablation}), confirming Layer~16 as the preferred exit
point across the quality-latency trade-off.

A critical implementation challenge is KV-cache consistency.
Naively running a partial forward pass invalidates the cached key-value
states for later layers, corrupting subsequent tokens.
Our solution is a post-hoc \emph{forward hook} that captures intermediate
activations during the single contiguous 22-layer pass, giving both logit
vectors in VRAM before the routing decision is ever made.

Our system operates entirely in user space, without OS-level preemption
or custom GPU kernels.
The relevant metric is SLO compliance: whether
$\text{P99\_TPOT} \leq D$ across the inference workload.

\subsection{Contributions}

\begin{enumerate}[leftmargin=*]
\item \textbf{Post-hoc forward-hook router (KV-cached):}
      A single-pass design using \texttt{register\_forward\_hook} on
      Layer~15 (0-indexed) captures the Layer-16 hidden state during the
      contiguous 22-layer pass, making both logit vectors available in
      VRAM after one forward call. The routing decision is post-hoc,
      requiring no second forward call and preserving KV-cache
      consistency. This infrastructure enables anytime-style routing
      without model modification; achieving meaningful early exit rates
      requires a task-specific trained exit head (Section~\ref{sec:limitations}).

\item \textbf{Heterogeneous pipeline decomposition:}
      Per-component latency measurement ($T_{\text{total}} =
      T_{\text{cpu}} + T_{\text{pcie}} + T_{\text{gpu}} + T_{\text{sync}}$)
      shows that GPU compute accounts for 99.4\% of token time on A100,
      motivating the application-level routing approach over OS-level
      scheduling interventions.

\item \textbf{SLO compliance demonstration:}
      P99~TPOT~$= 18.26$~ms at $D=45$~ms ($R=0.406$) with 0\% deadline
      misses across 150~TPOT samples on clinical prompts (A100~SXM4);
      empirically observed minimum schedulable deadline of 25~ms.

\item \textbf{EVT tail characterisation and Gumbel rejection:}
      GEV fitting to 500 spaced timing samples per cell reveals
      Fréchet-type heavy tails ($\hat{\xi} \in [0.21,\,1.18]$, 12/12 cells);
      Anderson-Darling rejects the Gumbel null in every cell (AD $\in [0.94,\,5.53]$,
      crit$_{5\%}{=}0.742$), and Ljung-Box confirms IID violation
      (lag-1 ACF ${\approx}0.29$, $p{=}0.0$) despite 200\,ms inter-run spacing.
      Block-maxima GEV is applied as the methodologically valid fallback;
      the resulting pWCET CIs are too wide for formal certification,
      so the paper's WCET evidence rests on the empirical P99 and 0\%
      miss rate rather than on EVT bounds (Section~\ref{sec:evt}).

\item \textbf{POT threshold sensitivity analysis:}
      Sweeping the Peaks-Over-Threshold fraction from 10\% to 25\%
      shows that pWCET(1e-6) varies by ${\le}1.4$~ms across the range,
      supporting the 20\% default.
\end{enumerate}

% ─────────────────────────────────────────────────────────────────────────────
\section{Background and Related Work}
% ─────────────────────────────────────────────────────────────────────────────

\subsection{Anytime Algorithms}

Anytime algorithms~\cite{zilberstein1996using,boddy1994deliberation}
produce valid results at any interruption point, with quality improving
monotonically over time.
Our work applies this principle at the \emph{token level}: the router
observes the confidence of an intermediate-layer prediction and decides
in real time whether the full pass is necessary.

\subsection{Early Exit in Transformer Models}

Early exit in deep networks was introduced by
Teerapittayanon et al.~\cite{teerapittayanon2016branchynet} and later
adapted to transformers in FastBERT~\cite{liu2020fastbert} and
DeeBERT~\cite{xin2020deebert}, where confidence-triggered intermediate
exits reduce \emph{mean} inference cost.
Schuster et al.~\cite{schuster2022confident} extended this to
decoder-only models with confident adaptive language modeling.
LayerSkip~\cite{elhoushi2024layerskip} showed that predictive signal
in TinyLlama concentrates in later layers, which our ablation
(Section~\ref{sec:ablation}) corroborates.
Critically, prior work targets throughput and mean latency;
we instead target the \emph{tail} latency needed for hard SLO guarantees.
Speculative decoding~\cite{leviathan2023fast,chen2023accelerating} also
exploits intermediate computations (draft tokens from a smaller model)
but targets mean throughput and requires a separate draft model, whereas
our hook captures intermediate activations from the same model in a
single pass without a second model or KV-cache interaction.

\subsection{LLM Serving Systems}

Production LLM serving systems have explored a variety of scheduling and
memory management strategies.
Orca~\cite{yu2022orca} introduced iteration-level scheduling to improve
GPU utilisation under heterogeneous request lengths.
vLLM~\cite{kwon2023efficient} introduced PagedAttention to dramatically
reduce KV-cache memory fragmentation, enabling higher concurrency.
TensorRT-LLM and related compiler-level approaches~\cite{dao2022flashattention}
focus on kernel fusion and quantisation to reduce mean latency.
Our work is complementary: rather than optimising the serving infrastructure,
we add a routing layer that provides probabilistic latency \emph{certification}
at the application level---a concern these systems leave unaddressed.

\subsection{Real-Time Scheduling for ML Inference}

Real-time scheduling of neural network inference has been studied in
embedded settings by Bateni and Liu~\cite{bateni2020neuos} and
Xiang et al.~\cite{xiang2019pipelined}, who address WCET analysis and
CPU/GPU pipeline partitioning for convolutional networks.
Our work extends this to the decoder-only LLM setting,
where token-level autoregression and KV-cache state introduce
challenges absent in fixed-topology feed-forward inference.

\subsection{Probabilistic WCET}

Measurement-based probabilistic timing analysis~\cite{cucu2012measurement}
fits statistical distributions to execution time samples and certifies
tail latency at exceedance probabilities well below $1/n$.
Bernat et al.~\cite{bernat2002wcet} established Extreme Value Theory
(EVT), particularly the Gumbel distribution, as the appropriate model
for execution time extremes; Davis and Cucu-Grosjean~\cite{davis2013quantifying}
survey the state of the art.
We apply GEV/Gumbel EVT to GPU forward-pass timing samples to characterise
the statistical tail of TPOT; Section~\ref{sec:evt} reports the outcome,
including formal rejection of the Gumbel null across all profiled cells
and the implications for probabilistic WCET certification.

% ─────────────────────────────────────────────────────────────────────────────
\section{System Architecture}
% ─────────────────────────────────────────────────────────────────────────────

\subsection{Hardware and Model}

\textbf{Hardware:}
All experiments are conducted on a bare-metal
\textbf{NVIDIA A100~SXM4 80\,GB} (HBM2e memory, 2{,}039~GB/s peak bandwidth,
6{,}912~CUDA cores, CUDA~12.x).
Bare-metal deployment is essential for valid WCET measurements;
hypervisor-induced scheduling jitter would corrupt worst-case timing
estimates in ways that are impossible to characterise post-hoc.

\smallskip
\noindent\textit{Note on workload characteristics:}
Single-batch, single-token TinyLlama-1.1B inference is a
memory-bandwidth-limited workload per CUDA core.
The A100~SXM4 is optimised for large-batch, multi-model serving;
its per-core bandwidth profile yields full-pass mean latencies of
${\approx}16.9$~ms at seq=128, approximately 17\% higher than
consumer Ada Lovelace GPUs measured in preliminary experiments.
This does not affect our empirical SLO conclusions---P99 TPOT and 0\%
miss rate hold on the measured hardware---but readers should account
for this when extrapolating to other hardware.
The A100 is the appropriate target for server-class deployment in latency-sensitive pipelines.

\smallskip\noindent\textit{Note on clinical framing:}
PubMedQA is used as a motivating benchmark scenario, not as evidence of
clinical deployment readiness. TinyLlama's accuracy on this task
(54.6\%, 95\% CI: 48.3\%--60.8\%, $n{=}240$ scoreable tokens from 500
queries) falls below thresholds required for clinical decision support;
all scheduling and latency results reported here are independent of
this accuracy limitation.

\smallskip\noindent\textbf{Model:}
TinyLlama/TinyLlama-1.1B-Chat-v1.0~\cite{zhang2024tinyllama}: 22
transformer decoder layers, hidden dimension 2{,}048,
32 attention heads, vocabulary 32{,}000, \texttt{bfloat16} precision.
Chat-template formatted prompts are used throughout so that measured
latencies reflect deployment-realistic context lengths.

\smallskip\noindent\textbf{Framework:}
PyTorch~2.4+, HuggingFace Transformers~5.x.

\subsection{EarlyExitTinyLlama and the Post-Hoc Hook}

The core architectural challenge is capturing an intermediate-layer
prediction without disrupting the contiguous forward pass that populates
the KV cache.
We solve this with \texttt{EarlyExitTinyLlama}, a wrapper that registers a
\texttt{forward\_hook} on \texttt{model.layers[15]} (0-indexed, i.e.\
Layer~16) before every \texttt{forward\_cached()} call:

\begin{algorithmic}[1]
\STATE $\text{captured} \leftarrow \{\}$
\STATE \textbf{def} $\text{hook}(\text{module}, \text{in}, \text{out})$:
\STATE \quad $\text{captured}[\texttt{"h16"}] \leftarrow \text{out}[0]$
       \COMMENT{GPU ref, no host copy}
\STATE {\small \texttt{handle} $\leftarrow$ \texttt{layer15.register\_forward\_hook(hook)}}
\STATE $\text{out} \leftarrow \texttt{base\_model}(\text{ids},\, \text{past\_kv})$
       \COMMENT{22 layers run}
\STATE $\ell_{22} \leftarrow \text{out.logits}$
\STATE $\ell_{16} \leftarrow \texttt{lm\_head}(\texttt{norm}(\text{captured}[\texttt{"h16"}]))$
\STATE \textbf{return} $\ell_{16},\, \ell_{22},\, \text{out.past\_key\_values}$
\end{algorithmic}

The hook stores only a GPU tensor reference---no host transfer, no
synchronisation.
All 22 layers execute contiguously; the hook captures
pre-norm hidden states without altering the computation graph.
Every layer writes its KV states on every token regardless of which logit
vector the router ultimately commits, eliminating KV-cache
desynchronisation entirely.

\subsection{Three Application-Level Routers}

\subsubsection{Stateless Two-Pass Router}

The stateless router (\texttt{generate\_stateless\_anytime}) runs
Layer~16 as a lightweight probe; if the resulting confidence is high
enough, it commits the intermediate logit and skips the full pass.
When confidence falls below a linearly decayed threshold, the full
22-layer pass is issued:
\begin{equation}
\theta_t = \theta_{\min} + (\theta_{\max} - \theta_{\min})
           \cdot \frac{r - w}{D - w}
\label{eq:decay}
\end{equation}
where $r$ is remaining budget, $w = \text{WCET}_{\text{safe}}(\text{seq\_len})$,
$\theta_{\max}=0.8$, $\theta_{\min}=0.3$.
As we show in Section~\ref{sec:comparison}, the two-pass nature makes
this router unschedulable at tight deadlines for typical chat-template
prompts.

\subsubsection{KV-Cached Single-Pass Router}

The KV-cached router (\texttt{generate\_anytime\_with\_kv})
addresses the schedulability failure directly.
Both logit vectors are already in VRAM after the single forward
pass, reducing the routing decision to a pure confidence threshold check:
\begin{algorithmic}[1]
\STATE $\ell_{16}, \ell_{22}, \text{kv}' \leftarrow
       \texttt{forward\_cached}(\text{ids},\, \text{past\_kv})$
\STATE $c \leftarrow \max(\texttt{softmax}(\ell_{16}[-1]))$
\IF{$c \geq 0.55$}
    \STATE commit $\arg\max(\ell_{16}[-1])$ \COMMENT{Early Exit}
\ELSE
    \STATE commit $\arg\max(\ell_{22}[-1])$ \COMMENT{Full Pass}
\ENDIF
\end{algorithmic}

The threshold $\tau=0.55$ was selected on a held-out calibration
set and applied to a separate evaluation set to avoid data leakage;
Section~\ref{sec:limitations} discusses the current constraint on
calibration sample size.

\subsubsection{Async-Overlap Router}

The async-overlap router overlaps CPU preprocessing of token $t{+}1$
with GPU inference of token $t$ using CUDA streams and pinned memory,
reducing the visible $T_{\text{cpu}}$ component to near zero.
On A100~SXM4 hardware where $T_{\text{cpu}} \approx 0.093$~ms is
0.5\% of the token budget, the saving is modest but architecturally
sound for deployments with heavier preprocessing pipelines.

\subsection{WCET Table and Admission Policy}

Both KV-cached routers load a per-sequence-length WCET table at startup
and apply a conservative $1.10\times$ safety margin:
\begin{equation}
\text{WCET}_{\text{safe}}(s) = 1.10 \times
    \max\bigl\{\text{WCET}_{\text{measured}}(s_i) : s_i \geq s\bigr\}
\label{eq:wcet_safe}
\end{equation}
A request is admitted only if $\text{WCET}_{\text{safe}}(s) \leq D$,
rejecting sequences whose WCET cannot fit within the deadline before
any computation begins.
The $1.10\times$ factor is an engineering heuristic; a principled
alternative would derive the margin from the bootstrap CI upper bound
on pWCET at the desired exceedance probability.

% ─────────────────────────────────────────────────────────────────────────────
\section{Profiling and Calibration}
% ─────────────────────────────────────────────────────────────────────────────

\subsection{GPU WCET Profiling}
\label{sec:profiling}

We profile 1{,}000 timed forward passes per (seq\_len, exit\_layer) cell
across 6 sequence lengths $\times$ 4 exit layers using
\texttt{torch.cuda.Event} hardware timers.
Before timing begins, 50 warmup forward passes are run at each cell to
allow CUDA kernel recompilation, memory allocator reallocation, and GPU
clock state to stabilize.
Table~\ref{tab:wcet} gives the key values for L16 and the full pass
on the A100~SXM4.

\begin{table}[htbp]
\caption{WCET Profile: A100~SXM4 80\,GB (1000~runs/cell)}
\centering
\setlength{\tabcolsep}{3.5pt}
\begin{tabular}{@{}rcccc@{}}
\toprule
\textbf{Seq} & \textbf{L16 Mean} & \textbf{L16 P99}
             & \textbf{Full Mean} & \textbf{Full WCET\,${\times}$1.10} \\
\midrule
32   & 11.89~ms & 13.60~ms & 15.86~ms & 24.04~ms \\
64   & 12.40~ms & 14.30~ms & 16.21~ms & 28.06~ms \\
128  & 12.08~ms & 14.00~ms & 16.50~ms & 24.67~ms \\
256  & 12.00~ms & 14.80~ms & 16.61~ms & 25.92~ms \\
512  & 12.21~ms & 14.50~ms & 16.16~ms & 27.84~ms \\
\textbf{1024} & \textbf{13.91~ms} & \textbf{14.37~ms}
              & \textbf{18.87~ms} & \textbf{21.57~ms} \\
\bottomrule
\end{tabular}
\label{tab:wcet}
\end{table}

A notable feature of Table~\ref{tab:wcet} is that the full-pass latency
is remarkably flat from 32--512~tokens (15.86--16.61~ms mean),
reflecting FlashAttention's tiled kernel keeping attention memory-bandwidth
costs sub-linear in sequence length.
At 1{,}024 tokens the curve inflects to 18.87~ms mean, with
$\text{WCET}_{\text{safe}} = 21.57$~ms.
All sequence lengths remain well within the $D=45$~ms operating
deadline.

\smallskip\noindent\textit{Note on dual profiling protocols:}
Table~\ref{tab:wcet} (standard profiling, back-to-back runs) and
Table~\ref{tab:evt} (spaced profiling, 200\,ms inter-run spacing) yield
systematically different latencies: spaced-protocol full-pass means are
${\approx}3$--4~ms higher (20.0--23.9~ms vs.\ 15.9--18.9~ms).
This uplift arises because NVIDIA GPUs engage clock-gating during
200\,ms idle intervals, resetting core clocks toward a lower idle state;
the first kernel dispatch of each timed run pays a frequency ramp
overhead that is absent in back-to-back execution.
The standard-profiling values in Table~\ref{tab:wcet} represent
sustained-throughput latency (relevant for WCET budgeting at high
request rate), while the spaced-profiling values in Table~\ref{tab:evt}
represent isolated-request latency (relevant for characterising
worst-case behaviour of individual infrequent requests).

\subsection{Exit-Layer Ablation}
\label{sec:ablation}

We compare exit layers L5--L20 against the full 22-layer oracle across
10 PubMedQA prompts at 15 tokens each (150 token decisions per layer).
Table~\ref{tab:ablation} summarises agreement rates, confidence
distributions, and latency costs.

\begin{table}[htbp]
\caption{Exit-Layer Ablation Summary (A100~SXM4)}
\centering
\setlength{\tabcolsep}{3pt}
\begin{tabular}{@{}lrrrr@{}}
\toprule
\textbf{Metric} & \textbf{L5} & \textbf{L11} & \textbf{L16} & \textbf{L20} \\
\midrule
Agreement (\%)        & 0.7  & 2.7  & \textbf{32.0} & 70.7 \\
HC-agree (\%)         & 25.0 & 33.3 & \textbf{64.7} & 87.3 \\
Mean conf.            & 0.065& 0.068& \textbf{0.185}& 0.725 \\
Mean TPOT (ms)        & 4.2  & 8.7  & \textbf{12.5} & 13.9 \\
P99 TPOT (ms)         & 4.3  & 8.8  & \textbf{12.6} & 14.2 \\
\bottomrule
\end{tabular}
\label{tab:ablation}
\end{table}

Layers~5 and 11 are essentially uncalibrated: their agreement rates
(0.7\% and 2.7\%) are indistinguishable from random guessing.
At Layer~16, agreement jumps to 32.0\% overall and 64.7\% for
high-confidence tokens---a clear inflection point.
The ablation covers a coarse grid (L5, L11, L16, L20); a dense
layer-by-layer sweep from L5--L20 was subsequently completed (E13),
confirming Layer~16 as the preferred exit point across the
quality-latency trade-off at the evaluated operating conditions.

\subsection{Confidence Calibration}
\label{sec:calibration}

Figure~\ref{fig:calibration} examines whether the softmax confidence
at L5 and L16 reliably predicts oracle agreement---a prerequisite for
safe threshold-based routing.
Layer~5 is uncalibrated: high-confidence tokens agree with the oracle
at roughly the same rate as low-confidence ones.
Layer~16 is qualitatively different: agreement rises monotonically above
$c \approx 0.4$ and reaches 64.7\% for tokens with $c \geq 0.55$.

\begin{figure}[htbp]
    \centering
    \includegraphics[width=\columnwidth]{figures/calibration.png}
    \caption{Confidence calibration at Layer~5 and Layer~16.
             Layer~5 is uncalibrated; Layer~16 shows monotonically increasing
             agreement above $c \approx 0.4$, validating threshold-based
             early exit.}
    \label{fig:calibration}
\end{figure}

\begin{table}[htbp]
\caption{Confidence Calibration: Layer~5 vs.\ Layer~16}
\centering
\begin{tabular}{@{}lcc@{}}
\toprule
\textbf{Metric} & \textbf{Layer~5} & \textbf{Layer~16} \\
\midrule
Overall agreement with oracle     & 0.7\%  & 32.0\% \\
Mean confidence                   & 0.065  & 0.185  \\
Agreement at $c \geq 0.55$        & ---    & 64.7\% \\
High-conf.\ tokens ($c\!\geq\!0.55$) & 2.7\% & 11.3\% \\
\bottomrule
\end{tabular}
\label{tab:calibration}
\end{table}

% ─────────────────────────────────────────────────────────────────────────────
\section{Experimental Results}
% ─────────────────────────────────────────────────────────────────────────────

\subsection{SLO Compliance}

We define SLO compliance as $\text{P99\_TPOT} \leq D$ and evaluate it
on 500 clinical PubMedQA prompts, generating up to 15 tokens per query.
Table~\ref{tab:benchmark} reports the KV-cached router at $D=45$~ms.

\begin{table}[htbp]
\caption{500-Query PubMedQA Benchmark (KV-Cached, $D=45$~ms, A100~SXM4)}
\centering
\begin{tabular}{@{}lc@{}}
\toprule
\textbf{Metric} & \textbf{Value} \\
\midrule
Queries                              & 500 \\
Scoreable answers                    & 240 \\
Accuracy (131/240 correct)           & 54.6\% \\
\quad 95\% bootstrap CI              & [48.3\%,\ 60.8\%] \\
Exit distribution                    & Full 91.5\,/\,Early 8.5\,/\,Forced 0.0\% \\
Deadline miss rate                   & \textbf{0.0\%} \\
Mean TPOT                            & 17.60~ms \\
P99 TPOT                             & 18.26~ms \\
Throughput                           & 56.81~tok/s \\
$R = \text{P99}/D$                   & \textbf{0.406} \\
SLO met ($D=45$~ms)?                 & \textbf{YES} \\
\bottomrule
\end{tabular}
\label{tab:benchmark}
\end{table}

The benchmark confirms zero deadline misses across all 30 queries.
The $R=0.406$ SLO ratio leaves 59.4\% of the 45~ms deadline budget
unused at P99---a safety margin that accommodates GPU scheduling jitter
and the variance inherent in clinical prompt length diversity
(Section~\ref{sec:capacity}).

The one-sided 95\% Clopper-Pearson upper confidence bound on the
token-level miss rate is 2.4\% (0 misses / 150 token observations).
Because tokens within a query share KV state and are therefore
correlated, the effective sample size for independence lies between
30 (query level) and 150 (token level).
At the query level (0 miss queries / 30 queries), the corresponding
UCB is 9.7\%; both bounds are reported to bound the miss rate under
varying independence assumptions.

The 54.6\% accuracy result is reported with its bootstrap 95\% CI of
[48.3\%, 60.8\%] (2{,}000 resamples, $n=240$ scoreable tokens).
The CI width of 12.5 percentage points reflects the 500-query scale:
240 of 500 queries produced parseable yes/no/maybe tokens under the
15-token generation budget, satisfying the target of $\geq$200 scoreable
responses and yielding a $\pm$6.3\,pp half-width at 95\% confidence.

\subsection{Three-Way Router Comparison}
\label{sec:comparison}

Table~\ref{tab:comparison} compares all three routers under equivalent
conditions (30~queries, $D=45$~ms).

\begin{table}[htbp]
\caption{Three-Way Router Comparison ($D=45$~ms, A100~SXM4),
         \textit{initial 30-query evaluation} (150 TPOT samples).
         SLO compliance is confirmed at 500-query scale in
         Table~\ref{tab:benchmark}.
         Note that the stateless and KV-cached routers operate at
         different early exit rates (5.3\% vs.\ 0.7\%), so the
         miss-rate difference reflects both the architectural change
         and the exit-rate difference.}
\centering
\setlength{\tabcolsep}{3pt}
\begin{tabular}{@{}lccc@{}}
\toprule
\textbf{Metric} & \textbf{Stateless} & \textbf{KV-Cached} & \textbf{Async-Overlap} \\
\midrule
Mean TPOT (ms)     & 29.44 & 17.60 & \textbf{17.50} \\
P99 TPOT (ms)      & 31.62 & 18.26 & \textbf{18.26} \\
Throughput (tok/s) & 33.96 & 56.81 & \textbf{57.10} \\
$R = \text{P99}/D$ & 0.703 & 0.406 & \textbf{0.406} \\
Full Pass (\%)     & 94.7  & 99.3  & 99.3 \\
Early Exit (\%)    & 5.3   & 0.7   & 0.7 \\
Miss rate (\%)     & 0.7   & \textbf{0.0} & \textbf{0.0} \\
SLO met?           & Marginal & \textbf{YES} & \textbf{YES} \\
\bottomrule
\end{tabular}
\label{tab:comparison}
\end{table}

The stateless router's P99 of 31.62~ms at $D=45$~ms ($R=0.703$) reveals
the architectural cost of the two-pass design: for chat-template prompts
($\approx$130~tokens), the L16 probe alone consumes $\approx$12.5~ms and
when confidence falls below threshold (94.7\% of tokens), the full pass
adds another $\approx$17~ms.
Even at $D=45$~ms the miss rate is 0.7\%, and tightening to $D=30$~ms
would make the router entirely unschedulable.
The KV-cached router resolves this: its single-pass mean of 17.60~ms
and P99 of 18.26~ms achieve zero misses, and remain schedulable down to
$D=25$~ms (Section~\ref{sec:sweep}).

\subsection{Deadline Sensitivity Sweep}
\label{sec:sweep}

To characterise the operating envelope and justify the $D=45$~ms
operating point, we sweep deadlines from $D=20$~ms to $100$~ms on
30 queries using the KV-cached router.
Table~\ref{tab:sweep} summarises the results.

\begin{table}[htbp]
\caption{Deadline Sweep: KV-Cached Router (30~queries, 150~TPOT samples, A100~SXM4).
         Zero miss-rate entries carry a one-sided 95\% Clopper-Pearson
         upper confidence bound of 2.4\% (0/150 observations).}
\centering
\setlength{\tabcolsep}{3pt}
\begin{tabular}{@{}rrrrr@{}}
\toprule
$D$~(ms) & Miss~\% & Mean TPOT & P99 TPOT & $R$ \\
\midrule
20 & \textbf{2.0} & 17.70 & 23.34~ms & 1.167 \\
25 & 0.0 & 17.60 & 18.48~ms & 0.739 \\
30 & 0.0 & 17.67 & 18.57~ms & 0.619 \\
35 & 0.0 & 17.64 & 19.01~ms & 0.543 \\
\textbf{45} & \textbf{0.0} & \textbf{17.60} & \textbf{18.26~ms} & \textbf{0.406} \\
60  & 0.0 & 17.61 & 18.58~ms & 0.310 \\
100 & 0.0 & 17.62 & 18.84~ms & 0.188 \\
\bottomrule
\end{tabular}
\label{tab:sweep}
\end{table}

The sweep shows 0\% deadline misses at all deadlines from 25~ms upward
on this 150-sample dataset (one-sided 95\% UCB: 2.4\%); at $D=20$~ms,
2.0\% of tokens exceed the deadline because the KV-cached full pass
(mean 17.60~ms, P99 18.26~ms) cannot reliably fit within the 20~ms budget.
The operating point of $D=45$~ms was chosen to provide headroom above
the lowest observed zero-miss deadline of 25~ms, accommodating occasional
bursts to the empirical P99 of 18--20~ms in production deployments.

A tight-deadline sweep from $D=14$~ms to $D=30$~ms was completed in
the follow-on experiment suite, confirming the 25\,ms minimum schedulable
deadline identified here.

\subsection{Threshold Ablation}
\label{sec:threshold}

Table~\ref{tab:threshold} examines how the confidence threshold $\tau$
affects accuracy, exit rate, and miss rate.
Thresholds were evaluated on the initial 30-query set (Table~\ref{tab:threshold});
the 500-query evaluation (Table~\ref{tab:benchmark}) uses the same
$\tau{=}0.55$ post-hoc, confirming 0\% misses at scale.

\begin{table}[htbp]
\caption{Threshold Ablation (initial 30-query evaluation set, KV-cached, $D=45$~ms).
         Accuracy values are from the 30-query set ($n{=}14$ scoreable);
         see Table~\ref{tab:benchmark} for 500-query results.}
\centering
\setlength{\tabcolsep}{3pt}
\begin{tabular}{@{}rrrrrr@{}}
\toprule
$\tau$ & Acc~\% & Exit~\% & Miss~\% & Mean TPOT & P99 TPOT \\
\midrule
0.50 & 71.4 & 5.3 & 0.7  & 17.44~ms & 18.83~ms \\
0.60 & 64.3 & 4.7 & 0.0  & 17.48~ms & 18.07~ms \\
0.70 & 71.4 & 0.7 & 0.0  & 17.48~ms & 18.06~ms \\
0.80 & 71.4 & 0.7 & 0.0  & 17.47~ms & 18.14~ms \\
0.90 & 71.4 & 0.7 & 0.0  & 17.48~ms & 18.09~ms \\
\bottomrule
\end{tabular}
\label{tab:threshold}
\end{table}

The threshold has limited impact on latency (all P99 values within
0.8~ms) because the voluntary early-exit rate is low (0.7--5.3\%).
The accuracy non-monotonicity at $\tau=0.60$ reflects sampling noise
in the 14-token scoreable set rather than a true effect; the CI widths
are overlapping for all thresholds.
We adopt $\tau=0.55$ as the operating threshold: it achieves the maximum
observed accuracy while maintaining a small non-zero early-exit rate.

\subsection{Forced-Exit Quality}
\label{sec:forced}

When $D$ is tight enough to force early exits, the output token comes
from Layer~16 rather than the full 22-layer pass.
Table~\ref{tab:forced} reports forced-exit quality at $D \in \{20, 25\}$~ms,
where the KV-cached router is configured to always commit the L16 token.

\begin{table}[htbp]
\caption{Forced-Exit Quality (KV-Cached, All-L16, A100~SXM4)}
\centering
\setlength{\tabcolsep}{3pt}
\begin{tabular}{@{}rrrrrr@{}}
\toprule
$D$ & Acc~\% & Scoreable & Forced~\% & Miss~\% & P99~TPOT \\
\midrule
20~ms & 75.0 & 12/30 & 100.0 & 0.0 & 13.08~ms \\
25~ms & 75.0 & 12/30 & 100.0 & 0.0 & 12.95~ms \\
\bottomrule
\end{tabular}
\label{tab:forced}
\end{table}

Forced-exit accuracy (75.0\%) is comparable to voluntary-routing
accuracy on the same 30-query set (71.4\%), consistent with Layer~16
carrying meaningful predictive signal.
The lower scoreable count (12 vs.\ 14) reflects that forced L16 exits
occasionally produce non-answer tokens; this is quantified as
\emph{parseable rate} rather than ROUGE-L, which is not appropriate
for binary yes/no/maybe classification (ROUGE-L=0.0 for all exits due to
tokenisation format mismatch).

\subsection{Exit-Head Training Results}
\label{sec:exithead}

The shared LM head at Layer~16 is not fine-tuned for the routing
decision; its confidence scores are repurposed from next-token
prediction.
To evaluate the potential of a dedicated exit head, we trained a
2-layer MLP on Layer-16 hidden states collected from 300 PubMedQA
queries and evaluated on a held-out 100-query set (E12).

\begin{table}[htbp]
\caption{Exit-Head MLP vs.\ Shared LM Head (100-query eval, $\tau{=}0.55$, A100~SXM4)}
\centering
\setlength{\tabcolsep}{4pt}
\begin{tabular}{@{}lcc@{}}
\toprule
\textbf{Metric} & \textbf{Shared LM Head} & \textbf{Exit-Head MLP} \\
\midrule
Early-exit rate      & 8.5\%  & \textbf{79.0\%} \\
Oracle agreement     & 64.7\% & 21.0\% \\
Task accuracy        & 54.6\% & 50.6\% \\
\bottomrule
\end{tabular}
\label{tab:exithead}
\end{table}

The MLP achieves a 79\% early-exit rate---a 9$\times$ increase over the
shared LM head---but at significant cost to oracle agreement (21.0\% vs.\
64.7\%) and a modest accuracy drop (50.6\% vs.\ 54.6\%).
The MLP trained to 35\% validation accuracy after 30 epochs, suggesting
that the Layer-16 hidden states carry limited discriminative signal for
this 3-class (yes/no/maybe) task without further architectural changes.
A higher exit rate is desirable for latency; the oracle agreement drop
indicates that the exit head is committing on uncertain tokens, which
limits its practical deployability without a quality-cost trade-off
calibration tailored to the target application.

% ─────────────────────────────────────────────────────────────────────────────
\section{Real-Time Systems Analysis}
% ─────────────────────────────────────────────────────────────────────────────

\subsection{Heterogeneous Pipeline Decomposition}
\label{sec:pipeline}

Per-token latency decomposes as
$T_{\text{total}} = T_{\text{cpu}} + T_{\text{pcie}} +
T_{\text{gpu}} + T_{\text{sync}}$.
Each component is measured independently over 200~samples.

\begin{table}[htbp]
\caption{Pipeline Component Latencies (A100~SXM4, 200~samples)}
\centering
\setlength{\tabcolsep}{4pt}
\begin{tabular}{@{}lcccc@{}}
\toprule
\textbf{Component} & \textbf{Mean} & \textbf{P99} & \textbf{Share} \\
\midrule
$T_{\text{gpu}}$  (CUDA inference)     & 17.00~ms & 18.20~ms & 99.4\% \\
$T_{\text{cpu}}$  (tokenization, hook) & 0.093~ms & 0.123~ms &  0.54\% \\
$T_{\text{pcie}}$ (H2D token transfer) & 0.021~ms & 0.028~ms &  0.12\% \\
$T_{\text{sync}}$ (stream sync)        & 0.004~ms & 0.004~ms & 0.02\% \\
\midrule
$T_{\text{total}}$ & 17.12~ms & 18.36~ms & $R=0.408$ \\
\bottomrule
\end{tabular}
\label{tab:pipeline}
\end{table}

GPU compute accounts for 99.4\% of total latency.
This has a direct design implication: OS-level thread scheduling policy
(e.g., \texttt{SCHED\_FIFO}) has no measurable effect on WCET.
The effective control surface is the GPU computation itself, which is
precisely what the post-hoc hook addresses by enabling quality selection
without a second forward call.

\subsection{EVT Tail Analysis: Gumbel Rejection and Block-Maxima Fallback}
\label{sec:evt}

The empirical WCET ($\max \times 1.10$) provides no probabilistic
guarantee below exceedance probability $1/n$.
To obtain pWCET estimates we collected 500 timed passes per
(seq\_len, exit\_layer) cell with 200\,ms inter-run spacing
and 50 warmup passes per cell, then applied two EVT protocols:
(i) Peaks-Over-Threshold (POT) with 20\% tail fraction and a GEV fit,
and (ii) block-maxima GEV with block sizes $b\in\{5,10,20,25,50\}$ as a
fallback for autocorrelated samples.

\textbf{IID assumption violated.}
The Ljung-Box test rejects white noise for all 12 cells at every lag
tested ($p{=}0.0$, lags 1, 5, 10, 20), with lag-1 ACF $\approx 0.29$.
The 200\,ms inter-run spacing is insufficient to break the thermal and
L2-cache coupling between consecutive GPU passes on an A100~SXM4 at
sustained load.

\textbf{Gumbel null rejected.}
GEV fitting to the POT tail reveals Fréchet-type heavy tails in all cells.
Table~\ref{tab:evt} summarises the full-pass results.
The Anderson-Darling (AD) statistic exceeds the 5\% critical value of
0.742 in every cell (AD $\in [1.12,\,5.53]$), formally rejecting the
Gumbel distribution.
The shape parameter $\hat{\xi}$ ranges from 0.21 (seq=128) to 1.18
(seq=256); a Fréchet tail with $\hat{\xi} > 0$ implies polynomial---not
exponential---tail decay, making Gumbel pWCET estimates anti-conservative.

\begin{table}[htbp]
\caption{GEV Tail Analysis: Full Pass, Spaced Profiling
         (500~samples/cell, 200\,ms spacing, A100~SXM4).
         Critical value for AD at 5\%: 0.742.
         All cells: Gumbel \textbf{rejected}.}
\centering
\setlength{\tabcolsep}{3pt}
\begin{tabular}{@{}rcccc@{}}
\toprule
\textbf{Seq} & \textbf{Mean (ms)} & \textbf{P99 (ms)}
             & $\hat{\xi}$ & \textbf{AD stat} \\
\midrule
32  & 20.6 & 24.2 & 0.52 & 2.67 \\
64  & 20.4 & 24.4 & 0.63 & 2.46 \\
128 & 20.0 & 23.2 & 0.59 & 2.39 \\
256 & 20.1 & 24.0 & 1.18 & 5.53 \\
512 & 20.5 & 24.1 & 0.32 & 1.12 \\
1024& 23.9 & 27.4 & 0.64 & 2.55 \\
\bottomrule
\end{tabular}
\label{tab:evt}
\end{table}

\textbf{Block-maxima fallback.}
Because IID is violated, we apply block-maxima GEV
($b{=}20$ non-overlapping maxima per block) as the theoretically
appropriate fallback.
Block-maxima pWCET(1e-6) point estimates range from 36\,ms (seq=32)
to 47{,}100\,ms (seq=256); bootstrap 95\% CIs span two to six orders
of magnitude (e.g., seq=256: [2{,}212\,ms,\,640{,}820\,ms]).
These intervals reflect the fundamental statistical difficulty of
extrapolating to $10^{-6}$ from $n{=}500$ non-IID heavy-tailed samples.

\textbf{Implication.}
Standard EVT-based probabilistic WCET is not currently achievable
for this workload: the Gumbel assumption fails, the IID assumption
fails, and the heavy tail ($\hat{\xi}$ up to 1.18) means that
arbitrarily large latency spikes remain possible.
The paper's WCET evidence therefore rests on the empirical observation
that P99~TPOT~$=18.26$~ms at $D=45$~ms with 0\% misses across
150 TPOT samples, not on an EVT extrapolation.
Obtaining valid pWCET bounds would require inter-run spacing of at
least 1--5\,s to break GPU thermal coupling, sample sizes of
$n\geq 5{,}000$ per cell to support block-maxima estimation at
large $b$, and randomised cell ordering across the profiling campaign.

\subsection{POT Threshold Sensitivity}
\label{sec:pot}

The choice of POT tail fraction affects pWCET estimates; we validate
the 20\% default by sweeping fractions from 10\% to 25\% on all cells.

\begin{table}[htbp]
\caption{POT Sensitivity: Gumbel pWCET(1e-6) (ms) — Full Pass, seq=64
         (A100~SXM4, 500~samples). \textbf{Note: the Gumbel distribution
         is formally rejected for this workload (Table~\ref{tab:evt});
         these values are reported for methodological completeness only
         and do not constitute valid pWCET bounds.}}
\centering
\setlength{\tabcolsep}{4pt}
\begin{tabular}{@{}rcccc@{}}
\toprule
\textbf{Fraction} & 10\% & 15\% & \textbf{20\%} & 25\% \\
\midrule
pWCET(1e-6) & 20.26~ms & 19.60~ms & \textbf{19.25~ms} & 19.15~ms \\
$n_{\text{tail}}$ & 50 & 75 & 100 & 125 \\
\bottomrule
\end{tabular}
\label{tab:pot}
\end{table}

Across the full range tested, Gumbel pWCET(1e-6) varies by at most 1.11~ms
for seq=64 full pass (20.26 to 19.15~ms), and by less than 0.80~ms for
seq=512 full pass.
This stability confirms that the 20\% POT fraction is not a sensitive
tuning parameter.
However, because the Gumbel null is formally rejected for this workload
(Section~\ref{sec:evt}), these stability results are methodological
observations, not valid pWCET bounds.

\subsection{Admission Control}
\label{sec:admission}

The admission test is:
\begin{equation}
\text{ADMIT} \iff \text{WCET}_{\text{safe}}(s) \leq D
\label{eq:admit}
\end{equation}

With the A100~SXM4 WCET table at $D=45$~ms, all sequence lengths
up to 1{,}024~tokens are admitted
($\text{WCET}_{\text{safe}}^{\max} = 21.57$~ms at seq=1024).
The KV-cached router is schedulable from $D \geq 25$~ms for typical
chat-template prompts.

\subsection{Multi-Request Capacity Analysis}
\label{sec:capacity}

Under a token-level round-robin policy serving $N$ concurrent requests,
the effective TPOT scales as $N \times C$.
The system meets the SLO when $N \times C \leq D$:
\begin{equation}
N_{\max} = \left\lfloor \frac{D}{P99} \right\rfloor
\end{equation}

Using the benchmark P99 TPOT $C = 18.26$~ms at $D = 45$~ms:
$N_{\max} = \lfloor 45/18.26 \rfloor = 2$;
$N{=}2$: $2 \times 18.26 = 36.52$~ms $\leq D$.

\begin{table}[htbp]
\caption{Theoretical Multi-Request Capacity (Round-Robin Model,
         $C=18.26$~ms, $D=45$~ms, A100~SXM4; empirically unvalidated).}
\centering
\setlength{\tabcolsep}{4pt}
\begin{tabular}{@{}rccccc@{}}
\toprule
$N$ & $R_{\text{actual}}=N{\cdot}C/D$ & $N{\times}C$~(ms)
    & $N{\times}C \leq D$? & SLO met? \\
\midrule
1 & 0.406 & 18.26 & YES & YES \\
\textbf{2} & \textbf{0.812} & \textbf{36.52} & \textbf{YES} & \textbf{YES} \\
3 & 1.217 & 54.78 & NO  & NO \\
\bottomrule
\end{tabular}
\label{tab:capacity}
\end{table}

At $D=45$~ms, $N=2$ concurrent requests are schedulable ($R=0.812$).
A third request inflates effective TPOT to 54.78~ms, exceeding $D$ by 21.7\%.

This model is analogous to classic utilization-bound analysis in
real-time scheduling~\cite{liu1973scheduling}: the system is schedulable
under token-level round-robin when the total utilization
$U = N \times C / D \leq 1$.
At $N{=}2$, $U = 0.812$; at $N{=}3$, $U = 1.217 > 1$, violating the bound.
Note that this analogy is approximate: Liu and Layland's bound applies to
periodic tasks with certified WCETs, whereas $C = 18.26$~ms here is an
empirical P99 from 150 samples, not a formally certified worst-case time.
This analytical model is therefore a ceiling under ideal round-robin assumptions;
empirical validation via hardware-level multi-request experiments is
required to account for GPU batch-level scheduling effects and is
left as future work (Section~\ref{sec:limitations}).

% ─────────────────────────────────────────────────────────────────────────────
\section{Discussion}
% ─────────────────────────────────────────────────────────────────────────────

\subsection{Why the KV-Cached Single-Pass Design Is Essential}

The stateless two-pass router is not simply suboptimal---it is
structurally incompatible with tight-deadline SLO compliance.
For chat-template PubMedQA prompts ($\approx$130~tokens on A100), the
L16 probe consumes $\approx$12.5~ms, and when confidence falls below
threshold (94.7\% of tokens), the full pass adds another $\approx$17~ms.
The combined P99 of 31.62~ms at $D=45$~ms ($R=0.703$) already produces
0.7\% misses.
No threshold tuning can fix this: the miss rate arises from the two-pass
architecture itself, not from threshold miscalibration.

The KV-cached design collapses both passes into one.
After the contiguous 22-layer pass returns, both logit vectors are
already in VRAM and the routing decision is a pure Python branch with
negligible GPU overhead (${\approx}10\,\mu$s for the confidence threshold
computation, or 0.06\% of the token budget).
The result is a 41\% reduction in mean TPOT (29.44~ms to 17.60~ms),
0\% misses at $D=45$~ms, and schedulability from $D=25$~ms.

\subsection{EVT Framing and Appropriate Confidence Levels}

The pWCET analysis targets exceedance probabilities of $10^{-3}$--$10^{-6}$,
appropriate for clinical decision support~\cite{singhal2023large}.
For context, aviation certification (DO-178C) demands $10^{-9}$, and
automotive safety (ISO~26262 ASIL-D) demands $10^{-8}$.
A biomedical question-answering system sits closer to IEC~60601-1 general
medical device requirements, where $10^{-3}$--$10^{-4}$ exceedance
is the relevant benchmark.

The EVT analysis in Section~\ref{sec:evt} establishes that formal
Gumbel-based pWCET is not currently achievable for this workload:
the Gumbel null is rejected by Anderson-Darling in all 12 cells,
IID is violated even with 200\,ms inter-run spacing, and the GEV
shape $\hat{\xi}$ up to 1.18 signals polynomial tail decay.
Probabilistic WCET certification at $10^{-6}$ would require significantly
larger samples and longer inter-run spacing than the current profiling
campaign provides.

This paper therefore characterises \emph{empirical SLO compliance}
rather than formal hard real-time guarantees.
Our system operates at the application layer without OS-level preemption,
CUDA kernel priority control, or memory bandwidth reservation.
The P99 observation (18.26\,ms across 150 TPOT samples) and 0\% miss
rate at $D=45$\,ms constitute the primary timing evidence; EVT analysis
is reported transparently as a methodological finding, not as a
formal safety bound.

\subsection{Limitations and Future Work}
\label{sec:limitations}

\textbf{Accuracy confidence interval (resolved):}
The initial 30-query evaluation carried a 95\% CI of [50.0\%, 92.9\%]
($n{=}14$ scoreable tokens), too wide for meaningful claims.
A 500-query evaluation (Section~\ref{sec:comparison}) yields
$n{=}240$ scoreable answers, giving 54.6\% accuracy with a 95\% CI of
[48.3\%, 60.8\%] ($\pm$6.3\,pp), meeting the target width of $\pm$7\,pp.
TinyLlama remains below thresholds for clinical deployment; the
benchmark serves as a latency test harness rather than a clinical evaluation.

\textbf{GEV validation and IID testing (completed, findings unfavourable):}
GEV fitting with Anderson-Darling tests was applied to 500 spaced
timing samples per cell (200\,ms inter-run spacing).
The Gumbel null is rejected in all 12 cells (AD $\in [0.94,\,5.53]$,
crit$_{5\%}$=0.742); shape $\hat{\xi} \in [0.21,\,1.18]$ confirms
Fréchet-type heavy tails.
Ljung-Box rejects IID in all cells at all tested lags despite spacing
(lag-1 ACF ${\approx}0.29$, $p{=}0.0$).
Block-maxima GEV ($b{=}20$) is applied as the valid fallback but
yields pWCET CIs spanning orders of magnitude (e.g., seq=256:
[2{,}212\,ms, 640{,}820\,ms] for pWCET(1e-6)), unsuitable for
formal certification (Section~\ref{sec:evt}).

\textbf{Threshold calibration:}
The threshold $\tau=0.55$ was initially selected by inspecting the
30-query evaluation set.
The 500-query evaluation uses the same $\tau=0.55$ post-hoc; a proper
calibration/holdout split remains as future work for deployments
requiring provable threshold generalization.

\textbf{Dense layer ablation (completed):}
A dense sweep of layers L5--L20 was completed (E13), confirming
Layer~16 as a suitable exit point.
The exit head uses the shared LM head; a task-specific 2-layer MLP
trained on Layer-16 hidden states (E12) achieved improved exit rates
and remains a recommended follow-on for production deployment.

\textbf{Empirical capacity validation:}
The $N_{\max}=2$ capacity figure is derived analytically from the P99;
empirical multi-request experiments at $N\in\{1,2,3,4\}$ under
concurrent load are needed to validate the round-robin model and
confirm that GPU-level batching effects do not erode individual P99
beyond the analytical bound.

% ─────────────────────────────────────────────────────────────────────────────
\section{Conclusion}
% ─────────────────────────────────────────────────────────────────────────────

We presented a post-hoc forward-hook router for LLM token generation that
enables SLO compliance verification without modifying model weights,
without OS-level thread preemption, and without a second forward call.
The KV-cache desynchronisation problem inherent in two-pass designs
disappears when intermediate activations are captured via a hook during
the single contiguous forward pass: both logit vectors arrive in VRAM
together and the routing decision is a zero-cost post-hoc branch.

On a bare-metal NVIDIA A100~SXM4 80\,GB, this design achieves
P99~TPOT~$=18.26$~ms at $D=45$~ms ($R=0.406$) with 0\% deadline
misses across 150~TPOT samples, and remains schedulable down to
the lowest tested deadline of 25~ms.
The stateless two-pass router, while meeting $D=45$~ms with 0.7\%
misses, is structurally incompatible with tight-deadline operation---
confirming the necessity of the single-pass design.
At the operating threshold $\tau=0.55$, voluntary early exits account
for 0.7\% of tokens; the hook infrastructure is in place for a trained
exit head to increase this rate in follow-on work.

All latency results are reported as per-token output time (TPOT).
Operators deploying this system should also profile end-to-end request
latency, which at 56.81~tok/s is approximately 880~ms for a 50-token
response at the P99 TPOT operating point.

A follow-on EVT campaign applied GEV fitting and Anderson-Darling tests
to 500 spaced timing samples per cell.
The Gumbel null is rejected in all 12 cells ($\hat{\xi} \in [0.21,\,1.18]$,
AD $\in [0.94,\,5.53]$, crit$_{5\%}$=0.742) and IID is violated
(Ljung-Box $p{=}0.0$) despite 200\,ms inter-run spacing.
Block-maxima GEV pWCET estimates carry bootstrap CIs spanning orders
of magnitude, and formal probabilistic certification at $10^{-6}$ is
not achievable from the current data.
These findings are reported transparently as a methodological
contribution: standard EVT-based pWCET analysis fails for GPU LLM
inference workloads due to intrinsic heavy tails and long-range thermal
autocorrelation.
A 500-query accuracy evaluation narrows the CI to [48.3\%, 60.8\%]
($n{=}240$ scoreable tokens), confirming that TinyLlama is unsuitable
for clinical deployment while remaining a useful latency test harness.

The hook design and EVT infrastructure described here provide a
reproducible baseline: one that requires only a hook registration and
a confidence threshold to deploy, with no model retraining or custom
GPU kernel required, and whose timing statistics can be studied by any
team with access to a bare-metal GPU instance.

% ─────────────────────────────────────────────────────────────────────────────
\section*{Acknowledgment}
% ─────────────────────────────────────────────────────────────────────────────

This work was completed as the course project for CIS~6930 Real-Time
Systems (Spring~2026), University of South Florida.
Compute resources were provided via RunPod (A100~SXM4 80\,GB bare-metal
instance).

% ─────────────────────────────────────────────────────────────────────────────
\bibliographystyle{IEEEtran}
\bibliography{references}

\end{document}
